\section{Supplementary Synthetic Bridge and Training Dynamics}
\label{sec:synthetic}

Synthetic models permit exact source alignment, dense checkpoints, and interventions that are expensive in LLMs. Their role is to develop measurements, expose confounds, and generate predictions. They do not establish that pretrained Qwen or Gemma models use the same circuit.

\subsection{What synthetic v10 contributes}

Synthetic v10 trains separate 4-layer, 4-head, \(d_{\rm model}=256\) causal transformers on counts 1--30 with learned absolute positions. Within that setting, it gives five concrete design lessons:
\begin{enumerate}
  \item final count is linearly readable in both modes, but the centroid path bends with depth, so local or piecewise steering is preferable to one global \(+1\) direction;
  \item Direct answer queries contain broad-source routing candidates, while trace queries contain matched-source bundles;
  \item progress/termination states and final-total states are causally dissociable;
  \item natural final-state transplants can transport donor count, whereas some trace-marker states control early close without transporting the final scalar;
  \item high trace-readout attention mass does not guarantee that the selected head output carries a total.
\end{enumerate}
These findings motivate the counter-curve, head-role, and mediation experiments above. They also justify the strict separation between representation, routing, and causal transport.

\begin{figure}[H]
\centering
\includegraphics[width=\textwidth]{synthetic_v10_attention_signatures.png}
\caption{Preliminary synthetic v10 signatures used to design the LLM study. Direct answer queries show broad source coverage; trace queries show matched-source retrieval; trace-to-answer readout is separate. The panels are descriptive until accompanied by query-local ablation and clean patch recovery.}
\label{fig:v10-attention}
\end{figure}

\subsection{Why v15 is not the final mechanism setting}

Version 15 uses Tiny Shakespeare backgrounds, \(d_{\rm model}=256\), RoPE or learned relative bias, all-sequence loss, and inserted marker targets. It changes several factors at once and is therefore a design stress test rather than a clean ablation. Final counts can approach ceiling even while whole-trace fidelity falls sharply, and position-only controls at trace/final sites can be highly predictive. The lesson is negative but important: correct final answers do not prove faithful iterative aggregation, and a strong progress probe may decode position or visible step.

We therefore retain v15 only for control development: index permutation, position matching, trace-fidelity decomposition, and separation of progress from final total. It is not used as the primary synthetic claim.

\subsection{v20 as the candidate training-dynamics setting}

Version 20 counts occurrences of a queried set of three natural characters in a 256-character Tiny Shakespeare window, restricted to totals 1--30. It uses query-first input, RoPE, atomic count tokens, paired data/order for separately trained Direct and trace models, and a two-stage objective: all-sequence loss through step 1500, then task-output loss through step 10,000. Both models have four layers, four heads, and \(d_{\rm model}=256\).

At the final one-seed checkpoint, the trace model reaches .912 exact count on held-out free-running examples, while the Direct model reaches .335. Whole-trace exactness and ordered-marker accuracy are also high in the trace model. The Direct output error does not exhibit a simple monotone increase with count in that sample, so the proposed count-state dispersion analysis remains necessary. Local repair artifacts suggest strong next-source and matched-source patching, including value-dominant transport, but the original run manifest marks part of the causal stage as failed. These results are hypotheses until the stage is rerun from a recorded revision with multiple seeds.

\begin{figure}[H]
\centering
\includegraphics[width=\textwidth]{synthetic_v20_mechanism_timeline.png}
\caption{Current one-seed v20 training timeline. Teacher-forced local abilities appear before free-running trace accuracy; next-source behavior appears before matched-source specialization. Because the objective changes at step 1500, this ordering is descriptive until replicated under no-switch and multi-seed controls. Labels are inherited from the internal report.}
\label{fig:v20-timeline}
\end{figure}

\subsection{Training-dynamics estimands}

At saved checkpoints we track free-running exact accuracy, teacher-forced local accuracy, the three routing scores in \cref{eq:broad-score,eq:target-score,eq:next-score}, count/progress SNR, query-local ablation damage, and clean patch recovery. Optimizer step is supplemented by semantic exposure: the cumulative number of supervised occurrences of each source identity, retrieval step, and answer count.

Three outcomes are distinguished:
\begin{enumerate}
  \item \textbf{smooth formation}: continuous mechanism metrics predict behavior without a discrete change point;
  \item \textbf{replicated transition}: a preregistered change-point model improves held-out likelihood and the ordering recurs across seeds and related causal metrics;
  \item \textbf{objective/exposure artifact}: the apparent transition tracks the loss switch or event frequency rather than a mechanism-specific phase.
\end{enumerate}
This follows work that uses mechanistic progress measures to explain apparently sudden learning~\citep{nanda2023grokking,olsson2022induction}.

\begin{todobox}
\textbf{S1---multi-seed mechanism dynamics.}
Train five paired initialization seeds (1234/2234/3234/4234/5234) with a seed-independent frozen data pool and paired Direct/CoT batch stream. Save 100-step snapshots. Select each final-checkpoint functional bundle once on a discovery split, freeze it, and track the same bundle backward; do not reselect the highest-scoring head at each checkpoint. At preregistered milestones run query-local ablation and clean patching. Report each seed's 10--90\% formation interval for next-source control, matched-source identity transport, progress update, final-total transport, and free-running exactness. Call an ordering replicated only if its sign agrees in at least four of five seeds and a paired interval excludes zero. Repair the failed causal stage and publish commands, manifests, and checksums.
\end{todobox}

\begin{todobox}
\textbf{S2---objective, exposure, and architecture controls.}
With at least three seeds per condition, vary one factor at a time: objective-switch step (0/1500/3000/10000), count sampling (natural/uniform), or context length (128/256/384). Plot formation against optimizer step, steps since the switch, and cumulative semantic exposure. Add a shared-backbone model with Direct/CoT mode tokens, loss-weight matching, digit-wise versus atomic count tokens, and held-out \(C>30\). If formation times track the switch or exposure, describe curriculum/exposure dynamics rather than an intrinsic phase transition.
\end{todobox}

\section{Threats to Causal Identification}

\paragraph{Retrieval versus aggregation failure.}
High list recall with low Direct count accuracy is compatible with aggregation or readout failure, but recall must be measured under the same prompt and decoding. Attention mass can dilute with length, and a multi-hop summary route can lack direct answer-to-source attention.

\paragraph{Visible-index and position leakage.}
Progress may be decoded from an item number, trace position, or answer distance. Primary anchors occur before source-bearing content or at item end, and analyses cross index permutation, index-free traces, fixed answer position, and variable trace extent.

\paragraph{Capability mismatch and floor effects.}
A clean CoT manifold and noisy Direct manifold may simply compare success with failure. We use all paired prompts, preregistered overlapping-difficulty cells, and four correctness strata. Common-correct analyses are explicitly selection-conditioned. If Gemma4-E4B Direct remains at floor, its role is failure-mechanism replication rather than unbiased mediation.

\paragraph{Different compute and decoding.}
Native CoT uses more generated tokens and may use a different sampling policy. Equal-length null scratchpads, fixed token budgets, deterministic matched decoding, and native-policy reruns separate semantic trace effects from extra compute and policy.

\paragraph{Post-treatment traces.}
Generated trace style is selected by the model. Canonical teacher forcing establishes local competence; free-running prefixes before the first error establish execution. Neither alone identifies a natural mediation effect.

\paragraph{Attention redundancy and superposition.}
One role may be distributed across heads, and one head may serve multiple roles. Bundle dose responses, union source coverage, query-local interventions, and clean recovery are preferred to naming one maximally scoring head.

\paragraph{Off-manifold interventions.}
Large centroid replacements can create unnatural states. Residual-preserving local paths and natural donor transplants are primary; activation norm, nearest-natural-state distance, and downstream KL are reported.

\paragraph{Cross-model non-homology.}
Qwen and Gemma may implement the same function with different depths and distributed components. Replication is defined by semantic endpoint and causal effect, not matching head indices.

\section{Related Work}

CoT prompting improves multi-step performance~\citep{wei2022cot}, but generated explanations can be unfaithful~\citep{turpin2023unfaithful}. Recent work studies information flow and local causal structure in reasoning traces~\citep{yang2026howcot,wei2026attricot}; counting adds exact source alignments and a known algorithmic target. Strong hidden-state probes can be diagnostic without being causal~\citep{yuan2026hidden}.

RULER extends single-needle retrieval to multi-needle and aggregation tasks~\citep{hsieh2024ruler}. Transformer counting can be limited by representation interference, vocabulary, and context length~\citep{yehudai2024count}. Our contribution is to compare the routing and state-update mechanisms of direct and trace-based counting in the same task.

The causal tools follow activation patching and causal tracing~\citep{meng2022rome}, circuit discovery~\citep{conmy2023acdc}, query-specific circuit analysis~\citep{wang2022ioi}, and activation steering~\citep{turner2023actadd}. We distinguish local ordered-token successor primitives~\citep{gould2023successor} from the next-source and continuation control required to cover a long-context source set.

\section{Limitations}

The causal LLM experiments in this draft are registered rather than complete. The existing behavioral audit mixes some checkpoint pairs, decoding policies, and protocol failure modes. Open-weight Qwen and Gemma results may not generalize to proprietary or non-autoregressive systems. Linear probes can miss nonlinear codes; flexible probes can memorize nuisances. Patching can be off-manifold, and ablation can trigger compensation. Native source alignment is imperfect after the first trace error. Synthetic models differ from LLMs in scale, pretraining, tokenization, and task familiarity. Finally, the response-surface forms were explored on existing data and must earn the name ``law'' by predicting the untouched Qwen3-14B cells.

\section{Reproducibility and Artifact Plan}

The final artifact will include request-level data or hashes, exact prompts, checkpoint/tokenizer revisions, chat templates, generation parameters, semantic-anchor parsers, intervention code, discovery/reporting splits, and seed-level tables. Each figure will be generated from a frozen table with a manifest and checksum. Numerical exactness, parse, format, and truncation remain separate. Synthetic checkpoints and optimizer states will be retained at the declared cadence.

\begin{todobox}
\textbf{R0---frozen analysis registry.}
Create one versioned YAML file containing: task unit; model/mode pairs; checkpoint hashes; core and mechanism cells; semantic sites; trace alignment rules; discovery/reporting groups; head scores and bundle selection; count-state and noise estimands; intervention dose grids; endpoint definitions; denominator gates; negative controls; multiplicity correction; minimum examples/seeds; exclusion codes; and the claim ladder in \cref{tab:claim-ladder}. Generate every table caption's sample size and provenance tag from this registry.
\end{todobox}
